\scanpage{95}
% Typescript inverted in the scan.
\mspage{-7-}
consequences of 1.9, of which only the following (cf. 1.7) will be stated
here, the paper comes to a close.

1.10 Theorem. \uline{Let} $\mathrm{k}$ \uline{be a perfect field of} $\mathrm{\dim\leq 1}$ \uline{and} $\mathrm{G}$ \uline{a
connected linear group over} $\mathrm{k}$. \uline{Then every conjugacy class over} $\mathrm{k}$ \uline{contains
an element over} $\mathrm{k}$.

After the remark that Kneser, using extensions of 1.8, has recently shown
(cf. 1.9) that $\mathrm{H^1(k,G)=0}$ if $\mathrm{k}$ is a $\mathrm{p}$-adic field and $\mathrm{G}$ a simply connected
semisimple group over $\mathrm{k}$, this introduction comes to a close.

\S2. \uline{Some recollections}. In this section we recall some known facts,
including some characterizations 2.11 of regular semisimple elements, and
establish some notations which are frequently used in the paper. If $\mathrm{k}$ is
a field, $\mathrm{k^*}$ is its multiplicative group. The term ``algebraic group'' is
often abbreviated to ``group''. If $\mathrm{G}$ is a group, $\mathrm{G_0}$ denotes its identity
component. If $\mathrm{x}$ is an element of $\mathrm{G}$, then $\mathrm{G_x}$ denotes the centralizer
of $\mathrm{x}$ in $\mathrm{G}$, and $\mathrm{x_s}$ and $\mathrm{x_u}$ denote the semisimple and unipotent parts
of $\mathrm{x}$ when $\mathrm{G}$ is linear. Assume now that $\mathrm{G}$ is a semisimple group, that
is, $\mathrm{G}$ is a connected linear group with no nontrivial connected solvable
normal subgroup. We write $\mathrm{r}$ for the rank of $\mathrm{G}$. Assume further that $\mathrm{T}$
is a maximal torus in $\mathrm{G}$ and that an ordering of the (discrete) character
group of $\mathrm{T}$ has been chosen. We write $\mathrm{\Sigma}$ for the system of roots relative
to $\mathrm{T}$ and $\mathrm{X_\alpha}$ for the subgroup corresponding to the root $\mathrm{\alpha}$.

2.1 $\mathrm{X_\alpha}$ \uline{is unipotent and isomorphic (as an algebraic group) to the}
% archive sheet number: 94

